\documentclass[11pt]{article}
\usepackage[T1]{fontenc}
\usepackage{lmodern,microtype,amsmath,amssymb,amsthm,mathtools,booktabs}
\usepackage[margin=1in]{geometry}
\usepackage[hidelinks]{hyperref}
\usepackage{xurl}
\newtheorem{theorem}{Theorem}[section]
\newtheorem{lemma}[theorem]{Lemma}
\newtheorem{corollary}[theorem]{Corollary}
\newtheorem{proposition}[theorem]{Proposition}
\theoremstyle{definition}
\newtheorem{definition}[theorem]{Definition}
\newtheorem{remark}[theorem]{Remark}
\newcommand{\Z}{\mathbb Z}
\newcommand{\F}{\mathbb F}
\newcommand{\R}{\mathrm R}
\newcommand{\Lturn}{\mathrm L}
\newcommand{\ind}{\mathbf 1}
\newcommand{\supp}{\operatorname{supp}}
\newcommand{\tape}{\mathcal C_n}
\newcommand{\eps}{\varepsilon}
\title{Maximal Separated Phase Clocks and Spatial Charges\\for Cyclic Langton Ants}
\author{Hao Ke}
\date{}
\begin{document}
\maketitle
\begin{abstract}
For a cyclic $n$-color Langton ant, we classify observables that separate the
complete coloring from the head pose and have a constant increment on every
update. The coloring statistic is allowed to be arbitrary, rather than assumed
additive, local, or spatially periodic. In one invariant heading--checkerboard
sector, the largest modulus of a unit-increment clock is $8n$ when the rule
contains odd numbers of both right and left turns, and $4n$ otherwise. Explicit
four-periodic spatial coefficients attain the bound. A coloring statistic
shared by both sectors has maximal modulus $4n$. For the classical two-color
ant this gives a sharp sixteen-phase clock, together with a complete
classification of separated cocycles over every abelian coefficient group.
For mixed rules we also classify all compactly supported coloring potentials:
they are unique finite sums of explicit order-two plaquette charges, with the
available face families determined by the two turn-count parities. The proofs
use forced additivity, local pose equations, and planar integration. Exact
Smith-form audits and independent literal implementations provide reproducible
controls. The resulting formulas describe the exact phase information available from
separated full-state observations.
\end{abstract}

\section{Introduction}
A cyclic Langton ant reads a cell, turns left or right according to its color,
advances that color through a finite cycle, and moves one lattice step.
The alternating horizontal and vertical moves partition its possible incoming
poses into two checkerboard sectors. This representation, and its interpretation
through local Truchet connections, go back to the study of generalized ants by
Gale, Propp, Sutherland and Troubetzkoy~\cite{gale}.

A natural way to seek algebraic information about the dynamics is to combine
a statistic of the entire coloring with a correction depending only on the
position and heading. Such observables include familiar spatial charges and
modular counters. We ask a completeness question: how much elapsed-time phase
can any observable of this separated form retain, even if its coloring
statistic is nonlinear and has unrestricted spatial dependence?

Our answer is exact. For a rule with $r$ right-turn colors and $l$ left-turn
colors, $r+l=n$, a sector-specific unit clock has maximal modulus
\begin{equation}\label{eq:mainM}
 M=\begin{cases}8n,&r,l\text{ both odd},\\4n,&\text{otherwise}.
 \end{cases}
\end{equation}
The classical rule has $M=16$. The extra factor of two requires spatial
resolution beyond period two, and disappears when one insists on a single
coloring statistic for both sectors. We additionally give a complete local
presentation for the separated cocycles of every mixed rule, and a basis for
all their compactly supported coloring potentials.

\paragraph{Relation to previous work.}
Jillhewar~\cite{jillhewar}, in the proposition labelled \texttt{prop:cocycle},
proves the classical additive-charge conditions $4w=0$ and the zero four-corner
sum, with sufficiency over $\Z/4\Z$. We use and explicitly recover this
zero-rate calculation; the familiar row and column charges are not newly
introduced here. McNamara~\cite{mcnamara} studies fixed finite ansatz families
with affine coordinate/time terms and zero-, one-, or two-bit site classes.
That study already explains why grouping the increments by turn removes the
order of the colors from the parameter equations. Our contribution is the
unrestricted spatial and arbitrary-coloring problem, its sharp nonzero phase
order, and the complete compact-charge basis. Order-blindness of the grouped
parameter equations is credited to that earlier observation, not claimed anew.

\section{All-state separated observables}
Let $\tau=(\tau_0,\ldots,\tau_{n-1})\in\{\R,\Lturn\}^n$ be nonempty. A
coloring $C:\Z^2\to\Z/n\Z$ belongs to $\tape$ if it is zero outside a finite
set. A state is $X=(C,p,h)$, where $p=(x,y)$ and $h\in\{0,1,2,3\}$ denotes
north, east, south, west, with step vectors
$e_0=(0,1),e_1=(1,0),e_2=(0,-1),e_3=(-1,0)$.
Set $r=\#\{i:\tau_i=\R\}$ and $l=\#\{i:\tau_i=\Lturn\}$,
where $0\le i<n$. Reading color $c=C(p)$ changes $h$ by $+1$ for $\R$ or
$-1$ for $\Lturn$, changes $c$ to $c+1\pmod n$, and moves in the new heading.
Denote this map by $T_\tau$.

The value
\[
 \eps=h-x-y\pmod2
\]
is conserved. We first fix $\eps=0$: even sites have incoming headings N/S
and odd sites E/W. The quarter-turn $(x,y,h)\mapsto(-y,x,h-1)$ carries the other
sector to this one and preserves the rule. All statements quantify over every
finite coloring and every pose in the stated sector, not merely over the orbit
of a blank input.

Let $G$ be an arbitrary abelian group and $\kappa\in G$. We call
\begin{equation}\label{eq:separated}
 I(C,p,h)=F(C)+\phi(p,h),\qquad
 I(T_\tau X)-I(X)=\kappa
\end{equation}
a separated cocycle of constant rate $\kappa$. No initial assumption is made
on $F$ except that it is a function on $\tape$. In particular it can initially
be nonlinear or depend on the entire finite coloring. A conserved quantity has
$\kappa=0$. A unit phase clock modulo $m$ has $G=\Z/m\Z$ and $\kappa=1$.

\begin{lemma}[Forced additivity]\label{lem:additive}
Every $F$ in \eqref{eq:separated} has a unique representation
\begin{equation}\label{eq:additive}
 F(C)=F(0)+\sum_{v\in\Z^2} f_v(c_v),\qquad f_v(0)=0.
\end{equation}
For a mixed rule, meaning $r,l>0$, there exist $a_v,b_v\in G$ such that
\begin{equation}\label{eq:increments}
 f_v(c)=R(c)a_v+L(c)b_v,\qquad r a_v+l b_v=0,
\end{equation}
For the representative $0\le c<n$, define
$R(c)=\#\{0\le i<c:\tau_i=\R\}$ and
$L(c)=\#\{0\le i<c:\tau_i=\Lturn\}$.
\end{lemma}
\begin{proof}
Choose an allowed heading at $v$ and a coloring with $C(v)=c$.
Write $C^{v,+}$ for the coloring equal to $C$ away from $v$ and with value
$c+1\pmod n$ at $v$; let $h'$ be the resulting heading. The update affects
no other cell. Equation~\eqref{eq:separated} gives
\[
 F(C^{v,+})-F(C)=\kappa+\phi(v,h)-\phi(v+e_{h'},h').
\]
The right side is independent of all other colors. Increment the finitely
many nonzero coordinates of $C$ from zero, in any fixed order. Their increments
therefore add, proving \eqref{eq:additive}. The single-cell colorings fix the
$f_v$ uniquely. Equal-turn colors have the same target pose and hence the same
increment: $a_v$ for R and $b_v$ for L. Summing these increments gives the first
part of \eqref{eq:increments}; going around the color cycle gives the second.
\end{proof}

The lemma does not say that every invariant of the ant is additive. The
separation in \eqref{eq:separated} is an essential hypothesis. A head correction
that also sees nearby colors is outside that hypothesis.

\section{A complete face presentation}
Write $d_v=a_v-b_v$. A square is an \emph{R-face} when its lower-left corner has
odd checker parity, and an \emph{L-face} when that corner has even parity.
Write $\partial S$ for the four corner vertices of a unit square $S$.
The directed four-pose circuit around an R-face uses four right turns; around
an L-face it uses four left turns. For example the L-face based at $(0,0)$ has
incoming poses $(0,0,\mathrm S),(1,0,\mathrm E),(1,1,\mathrm N),(0,1,\mathrm W)$.

\begin{theorem}[Face presentation]\label{thm:presentation}
For a mixed cyclic rule, weights \eqref{eq:increments} admit a pose correction
of constant rate $\kappa$ in sector zero if and only if
\begin{align}
 2d_v&=0,&r a_v+l b_v&=0 &&\text{at every site},\label{eq:site}\\
 \sum_{v\in\partial S}a_v&=4\kappa &&&&\text{for every R-face }S,\label{eq:Rface}\\
 \sum_{v\in\partial S}b_v&=4\kappa &&&&\text{for every L-face }S.\label{eq:Lface}
\end{align}
The pose correction is unique up to a constant.
\end{theorem}
\begin{proof}
At an even site, compare the N and S input equations with the same E or W
target pose. The two comparisons yield
\[
 \phi(v,\mathrm S)-\phi(v,\mathrm N)=a_v-b_v=b_v-a_v.
\]
Thus $2d_v=0$. At an odd site the identical calculation gives
$\phi(v,\mathrm W)-\phi(v,\mathrm E)=d_v$. The color-cycle relation was proved
in Lemma~\ref{lem:additive}. Summing the four pose equations on a face circuit
gives \eqref{eq:Rface} or \eqref{eq:Lface}.

For sufficiency let $U(v)$ be the correction at the primary heading, N on
even sites and E on odd sites. The opposite heading is assigned $U(v)+d_v$.
Orient horizontal grid edges from even sites to odd sites and vertical edges
from odd sites to even sites. Assign the following increments of $U$:
\begin{center}
\begin{tabular}{lll}
\toprule
Source parity & Edge to $u$ & $U(u)-U(v)$\\
\midrule
even & east & $\kappa-a_v$\\
even & west & $\kappa-a_v+d_v-d_u$\\
odd & north & $\kappa-a_v+d_v$\\
odd & south & $\kappa-a_v-d_u$\\
\bottomrule
\end{tabular}
\end{center}
Reverse traversal negates the increment. Its circulation is
$4\kappa-\sum a_v$ on an R-face, and $4\kappa-\sum b_v$ on an L-face;
the latter simplification uses $2d_v=0$. Both vanish.

Every finite closed walk in the square grid is a sum of elementary face
boundaries and edge backtracks. This can also be proved by successively
interchanging neighboring horizontal and vertical steps and removing
backtracks. Consequently path integration from the origin defines $U$ on all
of $\Z^2$. Substitute $\phi=U$ at primary headings and $\phi=U+d$ at opposite
headings into the local transitions. The table gives the required equations;
the remaining opposite-heading equalities use $2d=0$. The color-cycle relation
ensures that the functions $f_v$ wrap consistently. This proves sufficiency.
The edge gradients and the opposite-heading differences determine every pose
once $U(0,0)$ is fixed, proving uniqueness.
\end{proof}

This presentation includes arbitrary spatial coefficients and arbitrary
abelian targets. The rule order still enters the actual color functions through
$R(c),L(c)$; it disappears only from the equations on $a,b$. The latter
reduction is the order-blindness mechanism identified in~\cite{mcnamara}.

\section{The maximal phase period}
\begin{theorem}[Sharp phase order]\label{thm:period}
Let $M$ be given by \eqref{eq:mainM}. For any abelian group $G$, there is a
separated cocycle of rate $\kappa$ in one fixed sector if and only if
$M\kappa=0$. Its coloring coefficients and pose correction can be chosen
four-periodic in each coordinate. If the same coloring function must work in
both sectors, existence is equivalent to $4n\kappa=0$.
\end{theorem}
\begin{proof}
First suppose the rule is mixed. From \eqref{eq:site},
\[
 n a_v=l d_v,\qquad n b_v=-r d_v.
\]
Multiplying a face equation by $n$ gives
\begin{equation}\label{eq:necessaryorder}
 4n\kappa=l\sum_{\partial S_R}d_v
          =-r\sum_{\partial S_L}d_v,
\end{equation}
where any one face of each kind may be used. Since $2d_v=0$,
$8n\kappa=0$ is necessary. If $l$ is even the first expression vanishes; if
$r$ is even the second vanishes. This gives the required necessary order.

For existence when $4n\kappa=0$, write
$z(x,y)=\ind_{x\text{ odd},\,y\text{ odd}}$ and choose
\begin{equation}\label{eq:universalclock}
 a_v=b_v=4z(v)\kappa.
\end{equation}
Every face contains exactly one odd/odd vertex. All site and face equations
hold, so Theorem~\ref{thm:presentation} applies.

Now suppose $r,l$ are odd, so $n$ is even, and $8n\kappa=0$. Put
\begin{align}
 \chi(x,y)&=(-1)^{x+y},\notag\\
 u(x,y)&=(1-n/2)z(x,y)-(n/2)\chi(x,y)xy,\label{eq:u}\\
 a_v&=4u(v)\kappa,\qquad b_v=(1-n)a_v.\label{eq:doubled}
\end{align}
Then $d_v=na_v$ has $2d_v=0$, and
$r a_v+l b_v=n(1-l)a_v=0$, because $1-l$ is even.
The corner sum of $\chi xy$ is the checker sign of the lower-left corner.
It follows that the corner sum of $u$ is $1$ on R-faces and $1-n$ on L-faces.
Thus \eqref{eq:Rface} holds and
\[
 \sum_{\partial S_L}b_v=4(1-n)^2\kappa=4\kappa,
\]
since $4n(n-2)\kappa=0$. Theorem~\ref{thm:presentation} proves existence.
Taking $\kappa=1$ in $\Z/(8n)$ attains the upper bound.

The coefficients in \eqref{eq:doubled} have period four: the change after
$x\mapsto x+4$ is a multiple of $8n\kappa$, and similarly for $y$. The edge
increments in the proof of Theorem~\ref{thm:presentation} are also
four-periodic. Along either coordinate axis their four-edge sum is zero
(the terms alternate $\kappa,-\kappa$). Zero face circulation carries this
identity to every parallel line. Hence the integrated $U$, and then $\phi$,
can be four-periodic. Construction~\eqref{eq:universalclock} has the same
property.

For a constant all-R or all-L rule, the blank ant circles four cells. After
$4n$ updates its complete coloring and pose return, so $4n\kappa=0$ is
necessary for \emph{any} state observable of constant increment. For sufficiency set $a_v=b_v=4z(v)\kappa$, where $z(v)$ is the indicator
that both coordinates are odd. Then $d_v=0$, $na_v=0$, and every unit square
has corner sum $4\kappa$. The edge table in the sufficiency proof of
Theorem~\ref{thm:presentation} therefore has zero circulation even in this
constant-rule case. Its path integral $U$ supplies $\phi(p,h)=U(p)$, and direct
substitution verifies every allowed transition with
$f_v(c)=4c\,z(v)\kappa$. This constructs the clock whenever $4n\kappa=0$,
including a unit clock modulo $4n$.

Finally suppose one field function $F$ works in both sectors for a mixed
rule. Its increments $a_v,b_v$ are common to the sectors. Every geometric
face is an R-face in one sector and an L-face in the other. Therefore both
corner sums equal $4\kappa$, and every corner sum of $d$ vanishes. Multiplying
the sum of $a$ by $n$ gives $4n\kappa=0$. Construction
\eqref{eq:universalclock} uses one and the same $F$ in both sectors and is sharp.
\end{proof}

\begin{corollary}[Spatial resolution]\label{cor:resolution}
If the increments $a_v,b_v$ are two-periodic in either coordinate, a separated
constant-rate cocycle necessarily satisfies $4n\kappa=0$, even without a
periodicity assumption on $\phi$. When $r,l$ are odd, four-periodic coefficients
attain order $8n$.
\end{corollary}
\begin{proof}
Under two-periodicity in $x$, a four-corner sum is unchanged by moving the face
one column. The neighboring face has the opposite type. Hence both sums,
for $a$ and $b$, equal $4\kappa$ on every face. The final argument in
Theorem~\ref{thm:period} applies. The $y$ case is identical.
\end{proof}

Thus the small sublattice ansatz is a genuine restriction for detecting the
extra phase. For example RL has maximum $8$ under that restriction and $16$
without it; RRRL has $16$ and $32$, whereas RRLL has $16$ in both cases.

\section{The classical ant: a sixteen-phase formula}
For the rule RL, write $B$ for the complete finite black set. The two increments
are $a_v=w_v$, $b_v=-w_v$, and $4w_v=0$.

\begin{theorem}[Complete classical separated form]\label{thm:classical}
In sector zero, an observable \eqref{eq:separated} exists precisely when
$16\kappa=0$ and
\begin{equation}\label{eq:classicalform}
 F(B)=F(\varnothing)+\sum_{(x,y)\in B}
 (-1)^{x+y}\bigl(f(x)+g(y)-4xy\kappa\bigr),
\end{equation}
where $f,g$ take values in $G[4]=\{a:4a=0\}$. One may normalize $g(0)=0$;
then the functions are unique. The pose correction is unique up to a constant.
\end{theorem}
\begin{proof}
Lemma~\ref{lem:additive} first forces $F$ to be a sum of weights $w_v$.
For a square based at $(x,y)$ the two face equations reduce to
\[
 \sum_{v\in\partial S} w_v=-4(-1)^{x+y}\kappa.
\]
Let $A(x,y)=(-1)^{x+y}w(x,y)$. Its mixed difference is $-4\kappa$.
Telescoping over rectangles, including signed rectangles for negative
coordinates, gives
\[
 A(x,y)=A(x,0)+A(0,y)-A(0,0)-4xy\kappa.
\]
Take $f(x)=A(x,0)$ and $g(y)=A(0,y)-A(0,0)$. Both are 4-torsion. Applying
$4$ to the mixed difference forces $16\kappa=0$. Conversely, these conditions
make every weight 4-torsion and satisfy both face equations, so
Theorem~\ref{thm:presentation} applies.
\end{proof}

At zero rate this recovers the classical checker-signed row and column charges;
the underlying local cocycle condition is due to~\cite{jillhewar}. The theorem
adds the nonzero-rate term, sharp order, and the removal of an a priori
additivity assumption.

For an explicit unit clock modulo 16 take $f=g=0$. Set
\begin{equation}\label{eq:Phi}
 U(x,y)=\begin{cases}
 0,&x,y\text{ even},\\
 1,&x\text{ odd},\ y\text{ even},\\
 4y-2,&x,y\text{ odd},\\
 4x-1,&x\text{ even},\ y\text{ odd},
 \end{cases}\quad
 \Phi(x,y,h)=U(x,y)+8xy\ind_{h\in\{\mathrm S,\mathrm W\}}.
\end{equation}
Then
\begin{equation}\label{eq:Q16}
 Q(B,p,h)=-4\sum_{(x,y)\in B}(-1)^{x+y}xy+\Phi(p,h)\pmod {16}
\end{equation}
satisfies $Q(T_{\mathrm{RL}}X)=Q(X)+1$. This follows either by substitution
in the eight heading/color cases or from the integrated construction above.
For sector one evaluate the formula on the quarter-rotated state. In particular,
\begin{equation}\label{eq:time}
 t\equiv Q(X_t)-Q(X_0)\pmod {16}.
\end{equation}
All black cells enter the finite sum, including cells not yet read. The formula
is spatially four-periodic, and its evaluation needs no past trajectory.

\begin{corollary}
No separated unit clock modulo 32 exists for classical RL on one full sector,
even with an arbitrary nonlinear $F$ and arbitrary spatial pose correction.
For any mixed cyclic rule there is no nonconstant separated first integral
with values in a torsion-free abelian group, and no such clock of a fixed
positive real increment.
\end{corollary}
\begin{proof}
The first assertion is Theorem~\ref{thm:period}. For a torsion-free target,
that theorem forces $\kappa=0$. Equations~\eqref{eq:site} give $d=0$ and
$na=0$, hence $a=b=0$. The integrated pose correction is constant.
\end{proof}
This does not exclude strict clocks whose increment is positive but varies
with the state. It also does not put a bound of sixteen on information available
to a genuinely joint function of the complete state.

\section{All compact separated invariants}
We now set $\kappa=0$ and assume the rule is mixed. Call $F$ \emph{compact}
when it depends only on the colors in some fixed finite set of sites. By
Lemma~\ref{lem:additive}, this means that $a,b$ are finitely supported.

\begin{theorem}[Compact-charge basis]\label{thm:compact}
Modulo additive constants, the compact separated invariants over $G$ form
\begin{equation}\label{eq:compactmodule}
 \left(\bigoplus_{S\text{ an L-face}}G[2]\right)_{r\ \mathrm{even}}
 \ \oplus\
 \left(\bigoplus_{S\text{ an R-face}}G[2]\right)_{l\ \mathrm{even}},
\end{equation}
where a parenthesized summand is omitted when its subscript condition fails.
Every invariant is a unique finite sum of the plaquette generators below.
Thus classical RL has no nontrivial compact separated field invariant, over
any abelian group.
\end{theorem}
\begin{proof}
Since $2d=0$, we have $2a=2b$. Doubling the two face conditions shows that $2a$
has zero corner sum on every square. Therefore $\chi\,2a$ is a sum of a row
function and a column function. A finitely supported function of this form is
zero: outside its horizontal support the column function must be constant,
and outside its vertical support the row function must cancel that constant.
Hence $2a=2b=0$.

The color-cycle equation is now
$(r\bmod2)a+(l\bmod2)b=0$. If both counts are odd, $a=b$ and the all-face
zero-sum argument forces $a=b=0$. If $r$ is even and $l$ odd, $b=0$ and $a$
has only its R-face constraint. The other case is symmetric. If both counts
are even, the R-face constraint on $a$ and the L-face constraint on $b$ are
independent.

To solve the remaining R-face constraint, make a graph whose vertices are the
centers of the R-faces. Each original lattice vertex joins its two diagonally
opposite incident R-faces, and is an edge of this rotated square lattice. The
condition $\sum_{\partial S_R}a=0$ says exactly that this finitely supported
edge assignment is a cycle with coefficients in $G[2]$. Every such cycle is a
unique finite sum of elementary face boundaries. One elementary proof is to
integrate edge labels along dual paths starting outside a bounding box of
the support, assigning exterior faces coefficient zero. Zero boundary makes
the integration path independent; finite support leaves only finitely many
faces with nonzero coefficients.
If every edge difference is zero, all face coefficients equal the exterior
zero, proving uniqueness. Elementary faces of this graph correspond to the
original L-faces. This proves the first direct sum in
\eqref{eq:compactmodule}; exchange the face colors for the second.

It remains to provide the actual pose corrections. For an L-face $S$, let
$\Pi_L(S)$ be the four incoming poses of its all-L circuit. If $r$ is even,
and $z\in G[2]$, define
\begin{equation}\label{eq:compactgenerator}
 J^R_{S,z}(C,p,h)
 =z\left(\sum_{v\in\partial S}R(c_v)+\ind_{(p,h)\in\Pi_L(S)}\right).
\end{equation}
The wrap in color is well defined because $r$ is even. A right turn at a
corner changes the field sum by $z$ and moves the pose into or out of
$\Pi_L(S)$, changing its correction by $z$; all other transitions change
neither contribution. Thus the total is conserved. When $l$ is even the
analogous $J^L$ uses an R-face, $L(c)$ and the all-R pose circuit. Their
coefficients are exactly the elementary boundaries already classified, so
they realize and exhaust the direct sums, with unique normalized pose corrections.
\end{proof}

The condition is on the total counts, not on consecutive run lengths. For
example RLRL has both families in \eqref{eq:compactmodule}, despite its odd run
lengths. This is distinct from the even-run hypothesis used for recurrent
symmetry in~\cite{gale}; no conclusion about recurrence follows here.

\section{Verification and reproducibility}
The proofs above are analytic over all parameters and all abelian coefficient
groups. The accompanying implementation performs three kinds of exact control.
First, it checks every color transition at every four-periodic coordinate and
allowed heading for all rules of lengths one through eight, including color
wrap. It also evaluates local identities at very large integer coordinates.
Second, an independent raw linear presentation assigns a separate unknown to
every site/color and every pose, without assuming same-turn reduction. Exact
integer Smith forms on two- and four-periodic grids recover the respective phase
orders. Finite-box binary kernels and integer presentations recover the compact
basis and its order-two factors. These finite presentations do not replace the
all-plane proof.

Third, independent Python and C++ literal dynamics compare complete read-color
streams and complete final colorings for finite inputs in both sectors, including
unread remote cells. The Python head potential is computed by coordinate-path
integration; C++ reconstructs its potential independently on the raw pose graph.
Both periodically recompute the entire coloring contribution. The classical
closed formula~\eqref{eq:Q16} is checked separately. The deterministic report
records performed counts, negative controls and source hashes. Reproduction uses
\texttt{python phase\_cocycle/reproduce.py}; exact matrix checks require
SymPy and the native control requires a C++17 compiler.

\section{Limitations}
The maximality theorem concerns the separated, constant-rate, all-state class
\eqref{eq:separated}. It does not classify arbitrary joint-state functions,
pose corrections that inspect colors, variable-increment potentials, or
identities restricted to a selected orbit. A modular clock does not decrease
and cannot by itself be a finite-resource termination argument. Neither the
phase classification nor the compact-charge basis proves eventual periodicity,
period-104 uniqueness, the finite-support Highway Conjecture, or universality.
The new classification is not a claim that the earlier classical charges or
same-turn reduction were unknown. The verification is unformalized analytic
proof with exact computational controls, not a theorem-prover formalization.

\paragraph{Computational assistance.}
Mathematical exploration, source drafting and test implementation were assisted
by a language model. All claimed general results have explicit proofs above;
finite tests are reported only for their stated verification scope.

\begin{thebibliography}{9}
\bibitem{gale}
D. Gale, J. Propp, S. Sutherland and S. Troubetzkoy,
\emph{Further Travels with My Ant}, The Mathematical Intelligencer
\textbf{17}(3) (1995), 48--56.
\href{https://arxiv.org/abs/math/9501233v1}{arXiv:math/9501233v1}.

\bibitem{jillhewar}
A. Jillhewar, \emph{Finite-Support Periodic Highways of Langton's Ant:
Necessary Conditions, Transverse Exclusions, and Exact Search}, manuscript
source dated 17 August 2026. Additive-charge proposition
\texttt{prop:cocycle},
\href{https://github.com/Atharva12456/langtons-ant-highway/blob/76a7e234e577d992d086c204c3dea00ae72cf10a/paper/main.tex}{source revision 76a7e234e577}.

\bibitem{mcnamara}
T.S. McNamara, \emph{Invariant mining, stage 2: all moduli at once, sublattice
families, and what the rule can and cannot change}, research report,
22 October 2026. \texttt{turmite-autoresearch},
\href{https://github.com/timClicks/turmite-autoresearch/blob/8904956fbb5a0e3bce676a39e92f4543f3774c5f/experiments/archive/2026-10-invariant-mining-snf/README.md}{archived report, revision 8904956fbb5a}.
\end{thebibliography}
\end{document}
